\section*{Measurable Functions and Random Variables \hfill 4}

In measure-theoretic probability, probability mass function and probability density are not the most fundamental concepts. Instead, we introduce the notion of measurable functions to study random variables in a more general setting. By focusing on measurable functions, we can better understand how to combine two or more random variables into a new random variable. Thus, in this chapter, we primarily focus on measurable space, which consists of the sample space and a $\sigma$-algebra, without the probability measure.

The collection of random variables is closed under the usual arithmetic operations, meaning that the sum, difference, product, maximum, and minimum of two measurable functions are also measurable. Additionally, we can also show that the point-wise limit of an infinite sequence of measurable functions is measurable. These properties provide a great degree of flexibility for manipulating random variables.

\section*{4.1 Measurable Functions}

Consider a population of finite size and suppose we want to study the distribution of height and weight in this population. We randomly select a person from the population and measure his/her height and weight. The measurements depend on the sample we have selected and thus vary from sample to sample. We distinguish between the process of randomly sampling a person from the population and the procedure used to measure their height and weight. The former process is random and is described by a probability measure. Once a person is selected, there is no more randomness, and the height and weight are measured using a deterministic procedure.

The body mass index (BMI) is a measure of body fat based on a person's height and weight. It is calculated as the weight in kilograms divided by the height in meters squared and is a deterministic function. In other words, once we know a person's height and weight, we can compute their BMI without any uncertainty or variability.

However, if we are considering a population of individuals and sampling from that population, then the resulting BMI becomes a random variable. This is because the height and weight of each individual in the population may vary, and thus the resulting BMI values will also vary. But once we have selected a specific individual and measured their height and weight, the resulting BMI value is no longer random or uncertain. It becomes a fixed value that can be computed using the deterministic formula.

This example leads to the following definition of measurable function.

\begin{defbox}{4.1}
Let $(\Omega, \mathcal{F})$ be a measurable space. A function $X : \Omega \to \mathbb{R}$ is called a (real-valued) \textit{measurable function} if for any Borel set $B$ in $\mathcal{B}(\mathbb{R})$, we have
\[
X^{-1}(B) \triangleq \{\omega \in \Omega : X(\omega) \in B\} \in \mathcal{F}.
\]
\end{defbox}

We note that the definition of measurable function does not depend on any probability measure, only the $\sigma$-algebra is relevant. However, if a probability measure $\mu$ is given, and $(\Omega, \mathcal{F}, \mu)$ is a probability space, we refer to $X$ as a \textit{random variable}.

We will use the terms ``measurable functions'' and ``random variables'' interchangeably in subsequent discussions. The definition of random variables ensures that we can assign a probability to any set in the form
\[
\{\omega \in \Omega : X(\omega) \in B\},
\]
where $B$ is any Borel set. If we denote the probability measure by $P$, then the probability of the above event is written as
\[
P(\{\omega \in \Omega : X(\omega) \in B\}).
\]
As short-hand notation, we also write it as
\[
P(X(\omega) \in B) \text{ or } P(X \in B).
\]

In this chapter, we will study the properties of random variables that do not depend on how we assign probability to Borel sets. That is why we use the terminology ``measurable function''. We will include the assignment of probability later, after we have studied the class of functions to which we can assign probability.

In general, we can define a random variable that takes values in a set $\Omega'$ equipped with a $\sigma$-field $\mathcal{G}$. The range space $\Omega'$ can be Euclidean space $\mathbb{R}^d$, the set of complex numbers $\mathbb{C}$, or the set of infinite binary sequences. When the range is $\mathbb{R}^d$, the measurable function is called a random vector. Similarly, when the range is $\mathbb{C}$, the measurable function is a complex-valued random variable, and when the range is the set of infinite binary sequences, it is called a random sequence.

\begin{defbox}{4.2}
A function $f$ from measurable space $(\Omega, \mathcal{F})$ to measurable space $(\Omega', \mathcal{G})$ is called $(\mathcal{F}, \mathcal{G})$-\textit{measurable} if
\[
f^{-1}(B) \triangleq \{\omega \in \Omega : f(\omega) \in B\}
\]
is $\mathcal{F}$-measurable for all $B \in \mathcal{G}$. When the $\sigma$-field $\mathcal{G}$ in the codomain $\Omega'$ is understood from the context, we say that the function $f$ is $\mathcal{F}$-\textit{measurable}. If both $\mathcal{F}$ and $\mathcal{G}$ are understood, we will just say that the function $f$ is \textit{measurable}. When $\Omega$ is the real number line and $\mathcal{F}$ is the Borel algebra, we will refer to an $\mathcal{F}$-measurable function as a \textit{Borel measurable} function.
\end{defbox}

\noindent$\blacktriangleright$ \textbf{Notation} Although $f$ is a mapping from $\Omega$ to $\Omega'$, we sometime refer to $f$ as a map from the measurable space $(\Omega, \mathcal{F})$ to the measurable space $(\Omega', \mathcal{G})$, when we want to emphasize the associated $\sigma$-fields. In probability theory, random variables are generally denoted using capital letters, such as $X$ and $Y$. Technically speaking, random variables are functions and should be written as $X(\omega)$ and $Y(\omega)$. However, for the sake of convenience, we often use the simplified notation $X$ and $Y$.

\textbf{Example 4.1.1 (Measurable Functions)} \\
If $\mathcal{F}$ is the trivial $\sigma$-field $\{\emptyset, \Omega\}$, and $\mathcal{G}$ is the $\sigma$-field in which all singletons are $\mathcal{G}$-measurable, then an $(\mathcal{F}, \mathcal{G})$-measurable function is a constant function.

If $\mathcal{F}$ is the power set $\mathcal{P}(\Omega)$, then no matter what $\mathcal{G}$ is, all $(\mathcal{F}, \mathcal{G})$ functions from $\Omega$ to $\Omega'$ are measurable.

\textbf{Example 4.1.2 (Random Variable as Dictionary in Python)} \\
To illustrate the concept of a random variable when the sample space is a finite set, we can use a Python dictionary to implement a function that maps elements in a sample space to real numbers. Below is an example Python program.

\begin{verbatim}
from random import choice
Omega = ['a','b','c','d','e']      # sample space
X = {'a':1, 'b':2, 'c':3, 'd':4, 'e':5}  # random variable X
Y = {'a':0, 'b':3, 'c':2, 'd':-1, 'e':5} # random variable Y
omega = choice(Omega)  # pick a sample uniformly
print(f"X={X[omega]}, Y={Y[omega]}") # print the values of X and Y
\end{verbatim}

In this program, a random sample is generated using the choice function from the random library. We note that the two random variables are defined before any samples are drawn. Once a sample is generated, the values of $X$ and $Y$ are determined based on data stored in the dictionaries.

\noindent$\blacktriangleright$ \textbf{The Choice of $\sigma$-Fields in the Domain and Codomain} To maximize flexibility, we would like to choose the $\sigma$-field $\mathcal{F}$ in the domain to be as large as possible.

\begin{table}[h]
\centering
\caption{Codomains of measurable functions}
\label{tab:codomains}
\begin{tabular}{|l|l|l|}
\hline
$\Omega'$ & $\mathcal{G}$ & \\
\hline
$\mathbb{R}$ & $\mathcal{B}(\mathbb{R})$ & Real-valued random variable \\
$\bar{\mathbb{R}} = \mathbb{R} \cup \{\pm\infty\}$ & $\mathcal{B}(\bar{\mathbb{R}})$ & Real-valued random variable with infinity \\
$\mathbb{R}^d$ & $\mathcal{B}(\mathbb{R}^d)$ & $d$-dim. random vector \\
$\mathbb{C}$ & $\mathcal{B}(\mathbb{C})$ & Complex-valued random variable \\
$\{0, 1\}^\infty$ & $\mathcal{B}(\{0, 1\}^\infty)$ & Infinite random bit stream \\
$\mathbb{R}^\infty$ & $\mathcal{B}(\mathbb{R}^\infty)$ & Infinite random sequence \\
\hline
\end{tabular}
\end{table}

If we choose $\mathcal{F}$ to be the power set $\mathcal{P}(\Omega)$, then every function is measurable. On the other hand, we want $\mathcal{G}$ to be as small as possible, but it cannot be too small. It should contain all the basic events of interests. For example, when the codomain of the measurable function is $\mathbb{R}^n$, we usually take the Borel algebra on $\mathbb{R}^n$ as the $\sigma$-field in the codomain but do not necessarily consider complete measure space in the codomain.

By selecting different codomains and the associated $\sigma$-algebra, we obtain various types of random variables. Table 4.1 lists a few useful combinations.

We note that when the codomain of a measurable function is the extended real numbers $\bar{\mathbb{R}} \triangleq \mathbb{R} \cup \{\pm\infty\}$, the Borel algebra is generated by the ordinary open sets in $\mathbb{R}$ and the special open sets of the form $(a, \infty]$ and $[-\infty, a)$. Some textbooks refer to $\bar{\mathbb{R}}$-valued measurable functions as \textit{generalized random variables}.

\begin{thmbox}{4.1}
Let $(\Omega, \mathcal{F})$ be a measurable space and let $A$ be a subset of $\Omega$. The indicator function
\[
\mathbf{1}_A(\omega) \triangleq \begin{cases} 1 & \text{if } \omega \in A \\ 0 & \text{otherwise.} \end{cases}
\]
is $(\mathcal{F}, \mathcal{B}(\mathbb{R}))$-measurable if and only if $A$ is $\mathcal{F}$-measurable.
\end{thmbox}

\noindent\textit{Proof} We consider four cases. Suppose $B$ is a Borel set that contains 1 but not 0. Then the pre-image $\mathbf{1}_A^{-1}(B)$ is equal to the set $A$. If $B$ is a Borel set that contains 0 but not 1, then $\mathbf{1}_A^{-1}(B) = A^c$. When $B$ contains both 0 and 1, then $\mathbf{1}_A^{-1}(B) = \Omega$. Finally, when $B$ does not contain 0 nor 1, then $\mathbf{1}_A^{-1}(B) = \emptyset$. This exhausts all possibilities of Borel sets $B$ containing 0 or 1.

If $A$ is $\mathcal{F}$-measurable, then the four possible pre-images $A$, $A^c$, $\Omega$, and $\emptyset$ are all in $\mathcal{F}$, and hence, the indicator function $\mathbf{1}_A$ is $\mathcal{F}$-measurable. Otherwise, if $A$ is not in $\mathcal{F}$, then the indicator function is not $\mathcal{F}$-measurable. \hfill $\blacksquare$

For example, we can consider the set $\mathbb{Q}$ consisting of all rational numbers in the real number line. It is measurable because it is a countable union of singletons. The indicator function $\mathbf{1}_{\mathbb{Q}}$ is Borel measurable. However, if $C$ is the Cantor set, the indicator function $\mathbf{1}_C$ is also Borel measurable. However, if $V$ is the Vitali set constructed in Sect. 2.5, the indicator function $\mathbf{1}_V$ is not Borel measurable.